\documentclass{exam-zh}
\usepackage{edu-explanation}

\examsetup{
  question/show-answer = true,
  fillin/show-answer = true,
  solution/show-solution = show-stay,
}

\begin{document}

\eduSectionTitle{知识点：配方——把一般式化成顶点式}


\begin{eduSolutionBox}{核心公式：提取、半系数、合并}
  \begingroup
  \setlength{\parskip}{2.5mm}
当 $a\ne0$ 时，
\[
\begin{aligned}
y&=ax^2+bx+c\\
 &=a\left(x^2+\frac ba x\right)+c\\
 &=a\left[\left(x+\frac b{2a}\right)^2-\frac{b^2}{4a^2}\right]+c\\
 &=a\left(x+\frac b{2a}\right)^2-\frac{\Delta}{4a},
\qquad \Delta=b^2-4ac.
\end{aligned}
\]
\par
\textbf{第一步——提二次项系数：}$y=a\left(x^2+\dfrac ba x\right)+c$，使括号内 $x^2$ 的系数变为 $1$。\par
\textbf{第二步——按 $x^2+2mx$ 配方：}令 $2m=\dfrac ba$，则 $m=\dfrac b{2a}$，且 $x^2+2mx=(x+m)^2-m^2$。\par
\textbf{第三步——拆中括号：}$a[(x+m)^2-m^2]+c=a(x+m)^2-am^2+c$，再合并常数。\par
  \endgroup
\end{eduSolutionBox}




\begin{eduProblemBlock}[例题 1 · 首项系数为 1]
把 $y=x^2+2x+4$ 配成顶点式。
\end{eduProblemBlock}





\begin{eduAuxItem}{思路导航}
提二次项系数$\;\to\;$写成 $x^2+2mx$ 并配方$\;\to\;$拆中括号并合并\end{eduAuxItem}





\begin{eduDualExplain}
  \begin{eduExplainLeft}
    \eduColumnTitle{\eduIconBulb}{小贴士}
        \eduSideHint{先找 $2m$}{不要直接猜括号里的数；先写 $2m=2$，再得到 $m=1$。}
  \end{eduExplainLeft}

  \begin{eduExplainRight}
    \eduColumnTitle{\eduIconBook}{解答}
      \eduExplainStep{1}{提二次项系数}{$y=1(x^2+2x)+4$。}
      \eduExplainStep{2}{写成 $x^2+2mx$ 并配方}{$2m=2$，所以 $m=1$，$x^2+2x=(x+1)^2-1$。}
      \eduExplainStep{3}{拆中括号并合并}{$y=1[(x+1)^2-1]+4=(x+1)^2+3$。}
  \end{eduExplainRight}
\end{eduDualExplain}



\clearpage

\begin{eduProblemBlock}[例题 2 · 首项系数不为 1]
把 $y=2x^2+8x+5$ 配成顶点式。
\end{eduProblemBlock}





\begin{eduAuxItem}{思路导航}
提二次项系数$\;\to\;$写成 $x^2+2mx$ 并配方$\;\to\;$拆中括号并合并\end{eduAuxItem}





\begin{eduDualExplain}
  \begin{eduExplainLeft}
    \eduColumnTitle{\eduIconBulb}{小贴士}
        \eduSideHint{先提再找 $2m$}{先提取 $2$，括号内一次项系数才是 $4$；因此从 $2m=4$ 开始。}
  \end{eduExplainLeft}

  \begin{eduExplainRight}
    \eduColumnTitle{\eduIconBook}{解答}
      \eduExplainStep{1}{提二次项系数}{$y=2(x^2+4x)+5$。}
      \eduExplainStep{2}{写成 $x^2+2mx$ 并配方}{$2m=4$，所以 $m=2$，$x^2+4x=(x+2)^2-4$。}
      \eduExplainStep{3}{拆中括号并合并}{$y=2[(x+2)^2-4]+5=2(x+2)^2-3$。}
  \end{eduExplainRight}
\end{eduDualExplain}




\begin{eduProblemBlock}[例题 3 · $b/a$ 不是偶数]
把 $y=3x^2+5x-2$ 配成顶点式。
\end{eduProblemBlock}





\begin{eduAuxItem}{思路导航}
提二次项系数$\;\to\;$写成 $x^2+2mx$ 并配方$\;\to\;$拆中括号并合并\end{eduAuxItem}





\begin{eduDualExplain}
  \begin{eduExplainLeft}
    \eduColumnTitle{\eduIconBulb}{小贴士}
        \eduSideHint{分数仍然先写 $2m=$}{即使 $\dfrac53$ 不是偶数，也只需由 $2m=\dfrac53$ 得 $m=\dfrac56$。}
  \end{eduExplainLeft}

  \begin{eduExplainRight}
    \eduColumnTitle{\eduIconBook}{解答}
      \eduExplainStep{1}{提二次项系数}{$y=3\left(x^2+\dfrac53x\right)-2$。}
      \eduExplainStep{2}{写成 $x^2+2mx$ 并配方}{$2m=\dfrac53$，所以 $m=\dfrac56$，$x^2+\dfrac53x=\left(x+\dfrac56\right)^2-\dfrac{25}{36}$。}
      \eduExplainStep{3}{拆中括号并合并}{$y=3\left[\left(x+\dfrac56\right)^2-\dfrac{25}{36}\right]-2=3\left(x+\dfrac56\right)^2-\dfrac{49}{12}$。}
  \end{eduExplainRight}
\end{eduDualExplain}




\begin{eduProblemBlock}[例题 4 · $a,b$ 都是分数]
把 $y=\dfrac12x^2-\dfrac34x+1$ 配成顶点式。
\end{eduProblemBlock}





\begin{eduAuxItem}{思路导航}
提二次项系数$\;\to\;$写成 $x^2+2mx$ 并配方$\;\to\;$拆中括号并合并\end{eduAuxItem}





\begin{eduDualExplain}
  \begin{eduExplainLeft}
    \eduColumnTitle{\eduIconBulb}{小贴士}
        \eduSideHint{先算括号内系数}{提取 $\dfrac12$ 后，一次项系数是 $-\dfrac32$；再由 $2m=-\dfrac32$ 求 $m$。}
  \end{eduExplainLeft}

  \begin{eduExplainRight}
    \eduColumnTitle{\eduIconBook}{解答}
      \eduExplainStep{1}{提二次项系数}{$y=\dfrac12\left(x^2-\dfrac32x\right)+1$。}
      \eduExplainStep{2}{写成 $x^2+2mx$ 并配方}{$2m=-\dfrac32$，所以 $m=-\dfrac34$，$x^2-\dfrac32x=\left(x-\dfrac34\right)^2-\dfrac9{16}$。}
      \eduExplainStep{3}{拆中括号并合并}{$y=\dfrac12\left[\left(x-\dfrac34\right)^2-\dfrac9{16}\right]+1=\dfrac12\left(x-\dfrac34\right)^2+\dfrac{23}{32}$。}
  \end{eduExplainRight}
\end{eduDualExplain}




\begin{eduProblemBlock}[例题 5 · 一次项含根号]
把 $y=2x^2+2\sqrt3x-1$ 配成顶点式。
\end{eduProblemBlock}





\begin{eduAuxItem}{思路导航}
提二次项系数$\;\to\;$写成 $x^2+2mx$ 并配方$\;\to\;$拆中括号并合并\end{eduAuxItem}





\begin{eduDualExplain}
  \begin{eduExplainLeft}
    \eduColumnTitle{\eduIconBulb}{小贴士}
        \eduSideHint{根式也放进同一模型}{仍由 $2m=\sqrt3$ 求 $m$；只需注意 $m^2=\left(\dfrac{\sqrt3}{2}\right)^2=\dfrac34$。}
  \end{eduExplainLeft}

  \begin{eduExplainRight}
    \eduColumnTitle{\eduIconBook}{解答}
      \eduExplainStep{1}{提二次项系数}{$y=2(x^2+\sqrt3x)-1$。}
      \eduExplainStep{2}{写成 $x^2+2mx$ 并配方}{$2m=\sqrt3$，所以 $m=\dfrac{\sqrt3}{2}$，$x^2+\sqrt3x=\left(x+\dfrac{\sqrt3}{2}\right)^2-\dfrac34$。}
      \eduExplainStep{3}{拆中括号并合并}{$y=2\left[\left(x+\dfrac{\sqrt3}{2}\right)^2-\dfrac34\right]-1=2\left(x+\dfrac{\sqrt3}{2}\right)^2-\dfrac52$。}
  \end{eduExplainRight}
\end{eduDualExplain}




\end{document}
